% !PRINTED 161
\weberpage{161}
% [source-original mathematical slip, confirmed 2026-08-25 against scan p. 186 / printed p. 161]
% The print assigns degrees m and n, respectively, to f and phi, although printed p. 156
% defined f with degree n and phi with degree m. The product mn is unaffected. Preserved.
Wir können dann $f$ und $\varphi$ auch als homogene Functionen
vom Grade $m$ und $n$ der Veränderlichen $x, y_1, z_1$ ansehen, und
die Resultante wird dann eine homogene Function $mn^{\mathrm{ten}}$ Grades
$R(y_1, z_1)$ von $y_1$ und $z_1$. Setzen wir darin

\[
y_1 = b' + b\lambda, \qquad z_1 = c' + c\lambda,
\]

so erhalten wir eine Gleichung $mn^{\mathrm{ten}}$ Grades in $\lambda$, und wir
können, wenn $R(y_1, z_1)$ nicht identisch verschwindet, die Constanten
 $b, b', c, c'$ so annehmen, dass der Co\"efficient der höchsten
Potenz von $\lambda$, nämlich $R(b, c)$, von Null verschieden ist. Dann
entspricht jeder Wurzel der Resultante nur ein Werthsystem
des Verhältnisses $x:y:z$, und die Anzahl der gemeinschaftlichen
Wurzeln von $f$ und $\varphi$ ist also höchstens gleich $mn$. Um den Satz
aussprechen zu können, dass die Zahl der gemeinsamen Wurzeln
immer gleich $mn$ ist, muss man noch ein Uebereinkommen treffen,
wonach gewisse dieser Werthpaare mehrfach zu zählen sind.

Der Satz, der, in geometrischem Gewande, besagt, dass sich
zwei Curven $m^{\mathrm{ter}}$ und $n^{\mathrm{ter}}$ Ordnung in $mn$ Punkten schneiden,
wird das \gesperrt{Bezout'sche Theorem} genannt.

\begin{center}
§.\,50.

\gesperrt{Elimination aus mehreren Gleichungen.}
\end{center}

Die Principien der Elimination, die wir im Vorhergehenden
auf zwei Gleichungen angewandt haben, lassen sich ohne
Schwierigkeit auf mehrere Gleichungen mit einer entsprechenden
Anzahl von Unbekannten ausdehnen.

Nehmen wir zunächst die Gleichungen

\weq{1}{f(x) = 0, \quad \varphi(x) = 0, \quad \psi(x) = 0,}

der Grade $m, n, p$, deren Co\"efficienten $a_i, b_i, c_i$ zunächst als
unabhängige Veränderliche betrachtet werden sollen.

Wir wenden auf zwei von ihnen, etwa auf $f(x), \varphi(x)$, den
Algorithmus des grössten gemeinschaftlichen Theilers an und
erhalten einen letzten Rest, der eine rationale Function der $a$
und $b$ ist und der, von Nennern befreit, mit der Resultante $R$
von $f(x)$ und $\varphi(x)$ übereinstimmt, und einen vorletzten, der eine
% [corrected 2026-08-07, 8x zoom strip 8/8] "lineare Function" is letterspaced; "von $x$ ist."
% is set normally, so the run is the two words only. Claimed by the translation pass and
% independently by the sperrung pass on a different model.
\gesperrt{lineare Function} von $x$ ist. Die Bedingung für das Vorhandensein
 einer gemeinsamen Wurzel von $f = 0$ und $\varphi = 0$
ist das Verschwinden der Resultante, und der vorletzte Rest
%CONT>>
